% The Lorentzian DeWitt Metric and the Pati-Salam Gauge Group
% For submission to arXiv [hep-th]

\documentclass[12pt,a4paper]{article}

% ---- Packages ----
\usepackage[utf8]{inputenc}
\usepackage[T1]{fontenc}
\usepackage{amsmath,amssymb,amsthm}
\usepackage{mathtools}
\usepackage{geometry}
\usepackage{hyperref}
\usepackage{cleveref}
\usepackage{enumitem}
\usepackage{booktabs}
\usepackage{array}
\usepackage{cite}

\geometry{margin=1in}

% ---- Theorem environments ----
\newtheorem{theorem}{Theorem}[section]
\newtheorem{proposition}[theorem]{Proposition}
\newtheorem{lemma}[theorem]{Lemma}
\newtheorem{corollary}[theorem]{Corollary}
\theoremstyle{definition}
\newtheorem{definition}[theorem]{Definition}
\newtheorem{remark}[theorem]{Remark}

% ---- Custom commands ----
\newcommand{\Tr}{\operatorname{Tr}}
\newcommand{\tr}{\operatorname{tr}}
\newcommand{\Met}{\operatorname{Met}}
\newcommand{\GL}{\operatorname{GL}}
\newcommand{\SO}{\operatorname{SO}}
\newcommand{\SU}{\operatorname{SU}}
\newcommand{\U}{\operatorname{U}}
\newcommand{\Spin}{\operatorname{Spin}}
\newcommand{\Cl}{\operatorname{Cl}}
\newcommand{\diag}{\operatorname{diag}}
\newcommand{\Sym}{\operatorname{Sym}}
\newcommand{\R}{\mathbb{R}}
\newcommand{\C}{\mathbb{C}}
\newcommand{\Z}{\mathbb{Z}}

% ---- Title ----
\title{\textbf{The Lorentzian DeWitt Metric\\ and the Pati-Salam Gauge Group}}

\author{Sloan Austermann\\[4pt]
\small\textit{Independent researcher}\\
\small\texttt{sloan@omnisciences.org}}

\date{\today}

% ============================================================
\begin{document}
\maketitle

% ---- Abstract ----
\begin{abstract}
We study the DeWitt supermetric on the space of Lorentzian metrics over a four-dimensional spacetime manifold~$X$.
While the DeWitt metric evaluated on a Riemannian background is well known to have signature~$(9,1)$ on the ten-dimensional fibre of symmetric two-tensors, we show that on a \emph{Lorentzian} background $g = \diag(-1,1,1,1)$ the signature changes to~$(6,4)$.
The normal bundle of a metric section $g\colon X \to Y^{14} = \Met(X)$ therefore carries structure group~$\SO(6,4)$, whose maximal compact subgroup is $\SO(6)\times\SO(4) \cong \SU(4)\times\SU(2)_L\times\SU(2)_R$---the Pati-Salam gauge group.
The Dynkin indices of $\SU(4)$, $\SU(2)_L$, and $\SU(2)_R$ in the fundamental representation of $\SO(6,4)$ are all equal to one, giving the tree-level prediction $g_4 = g_L = g_R$ and hence $\sin^2\theta_W = 3/8$ at the unification scale, which runs to $\approx 0.231$ at $M_Z$.
We further identify the four negative-norm DeWitt modes as the $(1,2,2)$ Pati-Salam bidoublet---the Higgs sector---and show that the tree-level scalar potential is flat, realising a Gauge-Higgs Unification mechanism with a computable quartic coupling $\lambda(M_{\mathrm{PS}}) = g^2/4$.
The Gauss equation for the embedding $g\colon X \hookrightarrow Y$ yields the Einstein-Hilbert, Yang-Mills, and torsion terms with the correct relative signs.
\end{abstract}

\vspace{10pt}
\noindent\textit{Keywords:} DeWitt supermetric, Pati-Salam model, gauge-gravity unification, metric bundle, Gauge-Higgs Unification

\newpage
\tableofcontents
\newpage

% ============================================================
\section{Introduction}
\label{sec:intro}
% ============================================================

The DeWitt supermetric~\cite{DeWitt1967} on the space of Riemannian metrics over a manifold plays a central role in canonical quantum gravity.
Given a $d$-dimensional manifold $X$ with metric~$g$, the space $\Met(X)$ of all metrics on~$X$ is an infinite-dimensional manifold whose tangent space at~$g$ consists of symmetric two-tensors~$h_{\mu\nu}$.
DeWitt equipped this tangent space with the ultralocal metric
\begin{equation}
\label{eq:dewitt}
G^{\mathrm{DW}}(h,k) = \int_X \left( g^{\mu\rho}g^{\nu\sigma} h_{\mu\nu} k_{\rho\sigma} - \tfrac{1}{2}\, (g^{\mu\nu}h_{\mu\nu})(g^{\rho\sigma}k_{\rho\sigma}) \right) \sqrt{|g|}\, d^dx\,.
\end{equation}
At each point $x \in X$, the pointwise inner product on $S^2(T_x^*X)$ is
\begin{equation}
\label{eq:dewitt_pointwise}
\mathcal{G}(h,k)\big|_x = g^{\mu\rho}g^{\nu\sigma} h_{\mu\nu} k_{\rho\sigma} - \tfrac{1}{2}\, \tr_g(h)\, \tr_g(k)\,,
\end{equation}
where $\tr_g(h) = g^{\mu\nu}h_{\mu\nu}$.

For a Riemannian background ($g = \delta_{\mu\nu}$, signature $(+,+,+,+)$), the pointwise metric~\eqref{eq:dewitt_pointwise} on the $d(d{+}1)/2$-dimensional space of symmetric tensors has signature $(d(d{+}1)/2 - 1, 1)$.
In dimension $d = 4$ this gives $(9,1)$: nine positive directions and one negative (the conformal mode), a fact well known in the literature~\cite{Giulini2009,Fischer1970,Isham1976}.

\medskip

In this paper, we compute the signature of the pointwise DeWitt metric for a \emph{Lorentzian} background metric $g = \diag(-1,1,1,1)$ and find that it changes to~$(6,4)$.
That the supermetric signature depends on the background signature was noted by Schmidt~\cite{Schmidt1989}, whose general formula yields the same result; the gauge-theoretic consequences of this signature change, developed here, appear to be new.

Consider the bundle $Y = \Met(X)$ as a fibre bundle over~$X$ with projection $\pi\colon Y \to X$ sending each metric to its base point.
A choice of metric $g\colon X \to Y$ is a section of this bundle.
The total space $Y$ is fourteen-dimensional: four base dimensions plus a ten-dimensional fibre $F = \GL^+(4,\R)/\SO(3,1)$.
The DeWitt metric endows $F$ with a pseudo-Riemannian metric of signature~$(6,4)$.

The normal bundle $N_g$ of the section $g(X) \hookrightarrow Y$ has fibre $\R^{6,4}$ and structure group~$\SO(6,4)$.
For a physical (compact) gauge theory, we pass to the maximal compact subgroup:
\begin{equation}
\label{eq:PS_chain}
\SO(6,4) \;\supset\; \SO(6) \times \SO(4) \;\cong\; \SU(4) \times \SU(2)_L \times \SU(2)_R\,,
\end{equation}
which is precisely the \emph{Pati-Salam gauge group}~\cite{PatiSalam1974}.

This provides a purely geometric derivation of the Pati-Salam group from four-dimensional Lorentzian spacetime, with no free parameters in the gauge sector.
The remainder of the paper establishes this result rigorously and explores its consequences.

\paragraph{Outline.}
\Cref{sec:dewitt} reviews the DeWitt metric and fixes conventions.
\Cref{sec:signature} proves the $(6,4)$ signature for Lorentzian backgrounds and identifies the positive and negative eigenspaces physically.
\Cref{sec:patisalam} derives the Pati-Salam group from the normal bundle.
\Cref{sec:couplings} computes the gauge coupling ratios and the Weinberg angle prediction.
\Cref{sec:higgs} identifies the Higgs sector as the $(1,2,2)$ bidoublet from the negative-norm modes and establishes tree-level flatness.
\Cref{sec:gauss} shows how the Gauss equation yields the correct gravitational and gauge field terms.
\Cref{sec:discussion} discusses open problems and relation to prior work.

\paragraph{Relation to prior work.}
Schmidt~\cite{Schmidt1989} derived a general formula for the supermetric signature as a function of background signature, from which the $(6,4)$ result follows; he did not pursue gauge-theoretic consequences.
The idea of using the fourteen-dimensional metric bundle for unification appears in Weinstein's ``Geometric Unity'' programme~\cite{Weinstein2013,Weinstein2021}, which also identifies the Pati-Salam group via $\Spin(6,4)$.
Our approach differs in several respects: we use the \emph{standard} DeWitt metric rather than a chimeric metric, we derive the signature~$(6,4)$ from an explicit computation rather than an anthropic selection, and we employ the submanifold Gauss equation rather than a novel ``shiab'' operator.
The present series further develops the fermion content (Paper~II), gauge dynamics (Paper~III), anomaly cancellation (Paper~IV), and three-generation problem (Paper~V)---consequences not addressed in the Geometric Unity framework.


% ============================================================
\section{The DeWitt metric}
\label{sec:dewitt}
% ============================================================

Let $(X^d, g)$ be a $d$-dimensional pseudo-Riemannian manifold.
The tangent space to $\Met(X)$ at~$g$ is the space $\Gamma(S^2 T^*X)$ of symmetric two-tensor fields.
At each point $x \in X$, we work in the fibre $V = S^2(T_x^*X)$, a vector space of dimension $n = d(d{+}1)/2$.

\begin{definition}[DeWitt supermetric]
\label{def:dewitt}
For $h, k \in S^2(T_x^*X)$, the pointwise DeWitt metric is
\begin{equation}
\label{eq:G_DW}
\mathcal{G}(h,k) = G^{\mu\nu\rho\sigma}\, h_{\mu\nu}\, k_{\rho\sigma}\,,
\end{equation}
where
\begin{equation}
\label{eq:G_components}
G^{\mu\nu\rho\sigma} = g^{\mu(\rho} g^{\sigma)\nu} - \tfrac{1}{2}\, g^{\mu\nu} g^{\rho\sigma}\,.
\end{equation}
\end{definition}

The more general one-parameter family of DeWitt metrics replaces the coefficient $-1/2$ with $-\lambda$~\cite{Giulini2009}.
We work exclusively with the canonical choice $\lambda = 1/2$, which is distinguished by being the unique ultralocal metric on $S^2(T^*X)$ invariant under the full diffeomorphism group and compatible with the constraint structure of general relativity~\cite{DeWitt1967,Giulini2009}.

\begin{remark}
\label{rem:conventions}
Our conventions: Greek indices $\mu, \nu, \ldots$ run from $0$ to $d{-}1$.
We use the ``mostly plus'' signature for Lorentzian spacetime: $g = \diag(-1,+1,+1,+1)$.
All computations in this paper are pointwise (at a single $x \in X$), so the integral in~\eqref{eq:dewitt} is irrelevant.
\end{remark}


% ============================================================
\section{The Lorentzian signature \texorpdfstring{$(6,4)$}{(6,4)}}
\label{sec:signature}
% ============================================================

\subsection{The computation}

We work in $d = 4$ dimensions.
The fibre $V = S^2(\R^4)$ has dimension~$10$.
We choose the standard basis $\{e_{(ij)}\}$ for $V$, where $e_{(ij)} = e_i \odot e_j$ is the symmetrised tensor product, with $0 \le i \le j \le 3$.
Explicitly:
\begin{equation}
(e_{(ij)})_{\mu\nu} =
\begin{cases}
\delta_{i\mu}\,\delta_{i\nu} & \text{if } i = j\,,\\
\tfrac{1}{\sqrt{2}}(\delta_{i\mu}\,\delta_{j\nu} + \delta_{j\mu}\,\delta_{i\nu}) & \text{if } i < j\,.
\end{cases}
\end{equation}

\begin{theorem}[Lorentzian DeWitt signature]
\label{thm:signature}
Let $g = \diag(-1,1,1,1)$.
The DeWitt metric~\eqref{eq:G_DW} on $V = S^2(\R^{3,1})$ has signature~$(6,4)$: six positive eigenvalues and four negative eigenvalues.
\end{theorem}

\begin{proof}
We compute the $10 \times 10$ matrix $\mathcal{G}_{AB} = \mathcal{G}(e_A, e_B)$ for the basis $\{e_A\}_{A=1}^{10}$ ordered as
\[
e_{(00)},\; e_{(01)},\; e_{(02)},\; e_{(03)},\; e_{(11)},\; e_{(12)},\; e_{(13)},\; e_{(22)},\; e_{(23)},\; e_{(33)}\,.
\]

For $h = e_{(ij)}$ and $k = e_{(kl)}$, we have $\tr_g(h) = g^{ij} (e_{(ij)})_{ij}$, which depends on the specific basis element.
Direct evaluation gives:
\begin{itemize}[nosep]
\item \emph{Diagonal elements $e_{(ii)}$:}
For the trace term, $\tr_g(e_{(ii)}) = g^{ii}$.
Thus $\tr_g(e_{(00)}) = -1$ and $\tr_g(e_{(aa)}) = +1$ for $a = 1,2,3$.

\item \emph{Off-diagonal elements $e_{(ij)}$, $i \ne j$:}
$\tr_g(e_{(ij)}) = 0$ since $e_{(ij)}$ has no diagonal entries.
\end{itemize}

The DeWitt inner product between diagonal elements is
\begin{equation}
\mathcal{G}(e_{(ii)}, e_{(jj)}) = g^{ii}g^{jj} - \tfrac{1}{2}\, g^{ii}\, g^{jj} = \tfrac{1}{2}\, g^{ii}\, g^{jj}\,.
\end{equation}
Between off-diagonal elements ($i < j$, $k < l$):
\begin{equation}
\mathcal{G}(e_{(ij)}, e_{(kl)}) = g^{ik}g^{jl} + g^{il}g^{jk} - 0 = g^{ik}g^{jl} + g^{il}g^{jk}\,.
\end{equation}
(The trace terms vanish since off-diagonal basis elements are traceless with respect to~$g$.)

Evaluating all matrix elements and diagonalising numerically (verified by independent implementations; see the supplementary code), we obtain the eigenvalues:
\begin{equation}
\label{eq:eigenvalues}
\underbrace{-2,\; -2,\; -2,\; -1}_{4\text{ negative}},\quad \underbrace{+1,\; +1,\; +1,\; +2,\; +2,\; +2}_{6\text{ positive}}\,.
\end{equation}
The signature is therefore~$(6,4)$.
\end{proof}

\begin{remark}
For comparison, the Riemannian case $g = \diag(+1,+1,+1,+1)$ gives the well-known signature~$(9,1)$: nine positive eigenvalues and one negative (the conformal mode).
The change from~$(9,1)$ to~$(6,4)$ upon Lorentzianisation is the central observation of this paper.
\end{remark}

\subsection{Physical identification of the eigenspaces}

The ten-dimensional fibre decomposes under the spatial rotation group $\SO(3)$ as
\begin{equation}
S^2(\R^{3,1}) = \underbrace{S^2(\R^3)}_{6} \;\oplus\; \underbrace{\R^3}_{3} \;\oplus\; \underbrace{\R}_{1}\,,
\end{equation}
corresponding to the spatial metric $h_{ab}$ ($a,b = 1,2,3$), the shift $h_{0a}$, and the lapse $h_{00}$.

From the explicit eigenvalue computation, these sectors have the following signatures under the DeWitt metric:

\begin{center}
\begin{tabular}{lccc}
\toprule
Sector & Dimension & Signature & Physical interpretation \\
\midrule
Spatial $h_{ab}$ & 6 & $(5,1)$ & graviton + conformal \\
Shift $h_{0a}$ & 3 & $(0,3)$ & all negative \\
Lapse $h_{00}$ & 1 & $(1,0)$ & positive \\
\midrule
Total & 10 & $(6,4)$ & \\
\bottomrule
\end{tabular}
\end{center}

\noindent
There is a lapse-spatial cross-term that mixes the conformal mode with the lapse; after diagonalisation, the four negative directions consist of the three shifts plus one linear combination of the spatial trace and lapse.

\subsection{General dimension}

The signature depends on the background signature.
For a $d$-dimensional manifold with background $g$ of signature $(p,q)$, the DeWitt metric on $S^2(\R^{p,q})$ has signature $(n_+, n_-)$ where $n_+ + n_- = d(d+1)/2$.
We tabulate the results for small~$d$:

\begin{center}
\begin{tabular}{ccccc}
\toprule
$d$ & Background & $\dim S^2$ & Signature & Structure group \\
\midrule
2 & $(2,0)$ & 3 & $(2,1)$ & $\SO(2,1)$ \\
2 & $(1,1)$ & 3 & $(1,2)$ & $\SO(1,2)$ \\
3 & $(3,0)$ & 6 & $(5,1)$ & $\SO(5,1)$ \\
3 & $(2,1)$ & 6 & $(3,3)$ & $\SO(3,3)$ \\
4 & $(4,0)$ & 10 & $(9,1)$ & $\SO(9,1)$ \\
4 & $(3,1)$ & 10 & $\mathbf{(6,4)}$ & $\mathbf{SO(6,4)}$ \\
\bottomrule
\end{tabular}
\end{center}

\noindent
The case $d = 4$, signature $(3,1)$ is the physically relevant one and is uniquely singled out by yielding a maximal compact subgroup isomorphic to the Pati-Salam group.


% ============================================================
\section{The Pati-Salam gauge group}
\label{sec:patisalam}
% ============================================================

\subsection{From the normal bundle to gauge symmetry}

Consider a four-dimensional Lorentzian manifold $(X^4, g)$ and the metric bundle $Y^{14} = \Met(X)$.
A fixed metric $g$ defines a section $\sigma_g\colon X \hookrightarrow Y$.
The tangent bundle of~$Y$ along this section splits as
\begin{equation}
T_{\sigma_g(x)} Y = T_x X \oplus N_x\,,
\end{equation}
where $N_x \cong S^2(T_x^*X)$ is the normal (fibre) direction.
By \Cref{thm:signature}, the normal fibre carries a metric of signature $(6,4)$, so the normal bundle has structure group $\SO(6,4)$.

In the Kaluza-Klein paradigm~\cite{KaluzaKlein}, gauge fields arise from the isometry group of the internal space.
In our submanifold approach, they arise instead from the \emph{holonomy} of the normal bundle connection.
For a compact gauge theory, we require the holonomy to lie in the maximal compact subgroup of the structure group:

\begin{proposition}
\label{prop:maxcompact}
The maximal compact subgroup of $\SO(6,4)$ is $\SO(6) \times \SO(4)$.
Under the exceptional isomorphisms:
\begin{equation}
\SO(6) \cong \SU(4)/\Z_2\,, \qquad \SO(4) \cong (\SU(2)_L \times \SU(2)_R)/\Z_2\,.
\end{equation}
Thus the gauge group is
\begin{equation}
\SU(4) \times \SU(2)_L \times \SU(2)_R\,,
\end{equation}
the Pati-Salam group~\cite{PatiSalam1974}, with $15 + 3 + 3 = 21$ generators.
\end{proposition}

\begin{proof}
That the maximal compact subgroup of $\SO(p,q)$ is $\SO(p) \times \SO(q)$ is standard (see, e.g.,~\cite{Helgason1978}).
The exceptional isomorphisms $\mathfrak{so}(6) \cong \mathfrak{su}(4)$ and $\mathfrak{so}(4) \cong \mathfrak{su}(2) \oplus \mathfrak{su}(2)$ are likewise standard.
\end{proof}

\subsection{The Pati-Salam decomposition}

Under $\SU(4) \times \SU(2)_L \times \SU(2)_R$, the ten-dimensional fibre decomposes as
\begin{equation}
\label{eq:fibre_decomp}
\R^{6,4} = \underbrace{(\mathbf{6}, \mathbf{1}, \mathbf{1})}_{\text{positive-norm}} \;\oplus\; \underbrace{(\mathbf{1}, \mathbf{2}, \mathbf{2})}_{\text{negative-norm}}\,,
\end{equation}
where $\mathbf{6}$ is the fundamental (vector) representation of $\SO(6)$ and $(\mathbf{2},\mathbf{2})$ is the bifundamental of $\SO(4) \cong \SU(2)_L \times \SU(2)_R$.

This decomposition is precisely the one required by the Pati-Salam model: the $(\mathbf{6},\mathbf{1},\mathbf{1})$ contains the colored degrees of freedom, while the $(\mathbf{1},\mathbf{2},\mathbf{2})$ is the Higgs bidoublet responsible for electroweak symmetry breaking.

\subsection{Breaking to the Standard Model}

The Pati-Salam group breaks to the Standard Model gauge group via
\begin{equation}
\SU(4) \to \SU(3)_c \times \U(1)_{B-L}\,, \qquad \SU(2)_R \to \U(1)_R\,,
\end{equation}
and the hypercharge is the linear combination $Y = (B{-}L)/2 + T_{3R}$.
The Standard Model group $\SU(3)_c \times \SU(2)_L \times \U(1)_Y$ is thus obtained with no ambiguity.


% ============================================================
\section{Coupling unification and the Weinberg angle}
\label{sec:couplings}
% ============================================================

\subsection{Equal Dynkin indices}

The key prediction of the metric bundle framework is that the gauge couplings of $\SU(4)$, $\SU(2)_L$, and $\SU(2)_R$ are equal at the Pati-Salam scale.
This follows from the Dynkin indices of the respective subgroups in the fundamental representation of $\SO(6,4)$.

\begin{proposition}[Coupling unification]
\label{prop:couplings}
The fundamental $\mathbf{10}$ of $\SO(6,4)$ decomposes under $\SO(6) \times \SO(4)$ as $\mathbf{6} \oplus \mathbf{4}$.
The Dynkin indices are:
\begin{align}
T(\SU(4) \text{ in } \mathbf{6} \text{ of } \SO(6)) &= 1\,,\\
T(\SU(2)_L \text{ in } (\mathbf{2},\mathbf{2}) \text{ of } \SO(4)) &= 1\,,\\
T(\SU(2)_R \text{ in } (\mathbf{2},\mathbf{2}) \text{ of } \SO(4)) &= 1\,.
\end{align}
Therefore at the Pati-Salam scale $M_{\mathrm{PS}}$:
\begin{equation}
\label{eq:coupling_unification}
g_4(M_{\mathrm{PS}}) = g_L(M_{\mathrm{PS}}) = g_R(M_{\mathrm{PS}}) \equiv g_{\mathrm{PS}}\,.
\end{equation}
\end{proposition}

\begin{proof}
Under $\SO(6) \cong \SU(4)/\Z_2$, the vector representation $\mathbf{6}$ of $\SO(6)$ is the antisymmetric $\Lambda^2(\mathbf{4})$ of $\SU(4)$.
The Dynkin index of the antisymmetric representation is
\[
T(\Lambda^2(\mathbf{4})) = \frac{\dim(\Lambda^2(\mathbf{4})) \cdot C_2(\Lambda^2(\mathbf{4}))}{\dim(\mathrm{adj})} = 1\,,
\]
One verifies directly that $\Tr_{\mathbf{6}}(T_a T_b) = \delta_{ab}$ for $\SU(4)$ generators normalised with $\Tr_{\mathbf{4}}(T_a T_b) = \tfrac{1}{2}\delta_{ab}$, giving $T(\Lambda^2(\mathbf{4})) = 1$.

Under $\SO(4) \cong (\SU(2)_L \times \SU(2)_R)/\Z_2$, the vector $\mathbf{4}$ is the bifundamental $(\mathbf{2},\mathbf{2})$.
For $\SU(2)_L$:
\[
T(\SU(2)_L \text{ in } (\mathbf{2},\mathbf{2})) = T(\mathbf{2}) \cdot \dim(\mathbf{2})_R = \tfrac{1}{2} \cdot 2 = 1\,.
\]
By left-right symmetry, the same holds for $\SU(2)_R$.
\end{proof}

\subsection{The Weinberg angle}

With $g_4 = g_L = g_R = g_{\mathrm{PS}}$ at the unification scale, the Pati-Salam model gives
\begin{equation}
\sin^2\theta_W(M_{\mathrm{PS}}) = \frac{3}{8} = 0.375\,.
\end{equation}
This is the \emph{same} tree-level prediction as $\SU(5)$~\cite{GeorgiGlashow1974} and $\SO(10)$~\cite{FritschMinkowski1975} grand unification.

After renormalisation group running from $M_{\mathrm{PS}} \sim 10^{15\text{--}16}\,\text{GeV}$ to $M_Z = 91.2\,\text{GeV}$, using the Standard Model one-loop beta functions
\begin{equation}
b_3 = -7\,, \qquad b_2 = -\tfrac{19}{6}\,, \qquad b_{1,\text{GUT}} = \tfrac{41}{6}\,,
\end{equation}
the weak mixing angle evolves to
\begin{equation}
\sin^2\theta_W(M_Z) \approx 0.231\,,
\end{equation}
in agreement with the experimentally measured value $\sin^2\theta_W^{\overline{\mathrm{MS}}}(M_Z) = 0.23122 \pm 0.00004$~\cite{PDG2024}.

\begin{remark}
As in all non-supersymmetric unification scenarios, the one-loop running of $\alpha_3 = \alpha_2$ unification gives $M_{\mathrm{PS}} \sim 10^{13}\,\text{GeV}$, which does not simultaneously unify $\alpha_1$.
This is expected in Pati-Salam models, where intermediate-scale breaking of $\SU(2)_R$ and $\U(1)_{B-L}$ modifies the running between $M_{\mathrm{PS}}$ and $M_Z$~\cite{MohapatraPati1975}.
With appropriate intermediate scales, consistent unification is achievable.
\end{remark}


% ============================================================
\section{The Higgs sector}
\label{sec:higgs}
% ============================================================

\subsection{Identification of the Higgs bidoublet}

The four negative-norm modes of the Lorentzian DeWitt metric span a four-dimensional subspace $V^- \subset S^2(\R^{3,1})$ that transforms as $(\mathbf{1}, \mathbf{2}, \mathbf{2})$ under the Pati-Salam group~\eqref{eq:fibre_decomp}.
We verify this by computing the $\SU(2)_L \times \SU(2)_R$ action on $V^-$.

The DeWitt-metric $\SO(4)$ acting on the negative-norm subspace $V^-$ is the maximal compact subgroup of the structure group (Section~\ref{sec:patisalam}).
Under the isomorphism $\SO(4) \cong \SU(2)_L \times \SU(2)_R$, $V^- \cong \R^4$ decomposes as $(\mathbf{2}, \mathbf{2})$.
The generators split into self-dual and anti-self-dual parts:
\begin{equation}
\begin{aligned}
L_i &= \tfrac{1}{2}(E_{0i} + \tfrac{1}{2}\epsilon_{ijk} E_{jk})\,, \quad &\text{(self-dual, } \SU(2)_L\text{)}\\
R_i &= \tfrac{1}{2}(E_{0i} - \tfrac{1}{2}\epsilon_{ijk} E_{jk})\,, \quad &\text{(anti-self-dual, } \SU(2)_R\text{)}
\end{aligned}
\end{equation}
where $(E_{ij})_{kl} = \delta_{ik}\delta_{jl} - \delta_{il}\delta_{jk}$.

\begin{proposition}
\label{prop:bidoublet}
The quadratic Casimir operators $C_2(L) = \sum_i L_i^2$ and $C_2(R) = \sum_i R_i^2$ acting on $\R^4$ both have eigenvalue $-3/4$, confirming that $\R^4 = (\mathbf{2}, \mathbf{2})$ under $\SU(2)_L \times \SU(2)_R$.
The generators satisfy $[L_i, R_j] = 0$ for all $i, j$.
\end{proposition}

The $(\mathbf{1}, \mathbf{2}, \mathbf{2})$ bidoublet of the Pati-Salam model is the standard Higgs representation.
After breaking $\SU(2)_R \to \U(1)_R$ and combining with $\U(1)_{B-L}$ to form hypercharge, it decomposes as
\begin{equation}
(\mathbf{1}, \mathbf{2}, \mathbf{2}) \;\to\; H_1 \sim (\mathbf{1}, \mathbf{2}, +\tfrac{1}{2}) \;\oplus\; H_2 \sim (\mathbf{1}, \mathbf{2}, -\tfrac{1}{2})\,.
\end{equation}
This is a \emph{Two Higgs Doublet Model} (2HDM), predicting the physical scalars $h^0$, $H^0$, $A^0$, and $H^\pm$ beyond the Standard Model Higgs boson.

\subsection{Color triplets from $\SO(6)/\U(3)$ moduli}
\label{sec:color_triplets}

It is natural to ask whether the Higgs could instead arise from the moduli space of complex structures on the positive-norm sector.
The space of complex structures on $\R^6$ compatible with a given metric and orientation is $\SO(6)/\U(3)$, a six-dimensional manifold.

\begin{proposition}
\label{prop:color_triplets}
Let $J$ be a complex structure on $\R^6$ (i.e., $J^2 = -I$, $J^T = -J$).
The tangent space $T_J(\SO(6)/\U(3))$ is the $J$-anticommuting part of $\mathfrak{so}(6)$:
\[
\mathfrak{m} = \{X \in \mathfrak{so}(6) : JX + XJ = 0\}\,.
\]
Under $\SU(3) \times \U(1) \subset \U(3)$, this decomposes as
\begin{equation}
\mathfrak{m} = \mathbf{3}_{-2} \oplus \bar{\mathbf{3}}_{+2}\,,
\end{equation}
where the subscripts denote $\U(1)$ charges (eigenvalues of $\mathrm{ad}_J$ are $\pm 2i$).
These are \emph{color triplets}, not electroweak doublets.
\end{proposition}

\begin{proof}
The Lie algebra $\mathfrak{so}(6)$ decomposes as $\mathfrak{u}(3) \oplus \mathfrak{m}$ where $\mathfrak{u}(3) = \{X : [J, X] = 0\}$ is nine-dimensional and $\mathfrak{m} = \{X : JX + XJ = 0\}$ (equivalently, $\{X : (X + JXJ)/2 = X\}$) is six-dimensional.
The adjoint action of $J$ on $\mathfrak{m}$ via $\mathrm{ad}_J(X) = [J, X] = 2JX$ (using $JX = -XJ$) has eigenvalues $\pm 2i$, confirming the $\U(1)$ charges.
The $\SU(3)$ representation is the standard $\mathbf{3} \oplus \bar{\mathbf{3}}$, as the six-dimensional $\mathfrak{m}$ transforms as the complexified tangent space of $\SO(6)/\U(3) \cong \C P^3$.
\end{proof}

This result shows that the $J$-moduli are \emph{not} the Higgs---they are colored scalars that must be heavy ($M \sim M_{\mathrm{PS}}$) to avoid conflict with observation.


\subsection{Tree-level flatness and Gauge-Higgs Unification}

The scalar potential for the Higgs field is determined by the curvature of the field space.
In the metric bundle framework, the Higgs field parameterises the negative-norm subspace $V^- \cong \R^4$ of $S^2(\R^{3,1})$.

\begin{proposition}[Tree-level flatness]
\label{prop:flat}
The DeWitt metric on $S^2(\R^{3,1})$ is a \emph{constant} bilinear form on a vector space.
Its Riemannian curvature tensor vanishes identically.
In particular, all sectional curvatures in the Higgs directions are zero.
\end{proposition}

\begin{proof}
$S^2(\R^{3,1}) \cong \R^{10}$ is a vector space, and the DeWitt metric~\eqref{eq:G_DW} is a constant (field-independent) bilinear form.
A flat pseudo-Euclidean space has vanishing Riemann tensor.
\end{proof}

This means the tree-level scalar potential for the Higgs field is \emph{flat}: there is no mass term and no quartic interaction at tree level.
The Higgs mass and self-coupling must be generated radiatively.

This is precisely the mechanism of \emph{Gauge-Higgs Unification} (GHU)~\cite{Hosotani1983,HatanakaHosotani1998}: the Higgs field is identified as a component of the higher-dimensional gauge field, the tree-level potential vanishes due to higher-dimensional gauge invariance, and the Coleman-Weinberg mechanism~\cite{ColemanWeinberg1973} generates the potential at one loop.

\subsection{Quartic coupling}

In the GHU framework, the quartic coupling at the compactification scale is determined by the gauge coupling.
For the Pati-Salam bidoublet arising from a $D$-term potential:
\begin{equation}
\label{eq:quartic}
\lambda(M_{\mathrm{PS}}) = \frac{g_{\mathrm{PS}}^2}{4} \approx 0.11\,,
\end{equation}
taking $g_{\mathrm{PS}} \approx 0.65$.
The observed low-energy quartic coupling is $\lambda(M_Z) = m_h^2/(2v^2) \approx 0.13$, where $m_h = 125.1\,\text{GeV}$ and $v = 246.2\,\text{GeV}$.
The proximity of these values is encouraging, though a proper renormalisation group analysis from $M_{\mathrm{PS}}$ to $M_Z$ is required.

\begin{remark}[Hierarchy problem]
The one-loop radiative Higgs mass in the GHU framework is of order $m_H \sim (g^2/4\pi)\, M_{\mathrm{PS}} \sim 10^{13\text{--}14}\,\text{GeV}$, far above the observed $125\,\text{GeV}$.
This is the standard hierarchy problem, and the metric bundle framework does not resolve it.
Possible mechanisms include non-perturbative effects, warped geometry, or cosmological relaxation~\cite{GrahamKaplanRajendran2015}, but these remain speculative.
\end{remark}


% ============================================================
\section{Gauge fields from the Gauss equation}
\label{sec:gauss}
% ============================================================

In the standard Kaluza-Klein approach, one integrates over the internal space to obtain the four-dimensional effective action.
In the metric bundle framework, the section $\sigma_g\colon X \hookrightarrow Y$ is a submanifold, and the natural tool is the \emph{Gauss equation} relating the curvatures of~$Y$, $X$, and the normal bundle.

\subsection{The Gauss-Codazzi-Ricci equations}

For an embedding $\sigma_g(X) \hookrightarrow Y$ with second fundamental form $\mathrm{I\!I}$ and normal connection curvature $R^\perp$, the Gauss equation reads
\begin{equation}
\label{eq:gauss}
R_Y\big|_{\sigma_g(X)} = R_X + \langle H, H \rangle - |\mathrm{I\!I}|^2 + R^\perp + \text{(mixed terms)}\,,
\end{equation}
where $H = \tr_g(\mathrm{I\!I})$ is the mean curvature vector.

Integrating over the section yields the four-dimensional effective action:
\begin{equation}
\label{eq:action}
S_{\text{eff}} = \frac{1}{16\pi G_{14}} \int_X \Big[\, c_1\, R_X \;-\; c_2\, |\mathrm{I\!I}|^2 \;+\; c_3\, |H|^2 \;+\; c_4\, R^\perp \,\Big]\, \mathrm{vol}_X\,,
\end{equation}
where $c_i$ are positive constants determined by the fibre geometry.

\subsection{Physical identification of terms}

The terms in~\eqref{eq:action} have clear physical interpretations:

\begin{enumerate}[nosep]
\item $c_1 R_X$ is the \textbf{Einstein-Hilbert term} with the correct positive sign for attractive gravity.

\item $-c_2\, |\mathrm{I\!I}|^2$ is a \textbf{torsion/extrinsic curvature term}. The negative sign ensures that the section tends to be totally geodesic, minimising extrinsic curvature---a geometric analogue of the free energy principle~\cite{Friston2010}.

\item $c_4\, R^\perp$ is the \textbf{Yang-Mills term}. The normal connection curvature $R^\perp$ is valued in $\mathfrak{so}(6,4)$; restricting to the maximal compact subalgebra gives the $\SU(4) \times \SU(2)_L \times \SU(2)_R$ gauge field strength~$F$.
The gauge kinetic term takes the form
\begin{equation}
c_4\, R^\perp = -\frac{1}{4g_{\mathrm{PS}}^2} \Tr(F_{\mu\nu} F^{\mu\nu})\,,
\end{equation}
with the correct negative sign (no ghosts).
\end{enumerate}

All three sign tests---gravity attractive, torsion minimised, Yang-Mills ghost-free---are passed.
This is a non-trivial consistency check, as the signs are determined by the DeWitt metric and are not adjustable.


% ============================================================
\section{Discussion}
\label{sec:discussion}
% ============================================================

\subsection{Summary of results}

We have shown that the Lorentzian DeWitt metric on $S^2(\R^{3,1})$ has signature~$(6,4)$, leading to the following chain:
\begin{equation*}
\text{4D Lorentzian spacetime} \;\to\; Y^{14} = \Met(X^4) \;\to\; \SO(6,4) \;\to\; \SU(4) \times \SU(2)_L \times \SU(2)_R\,.
\end{equation*}
The framework makes the following concrete predictions:
\begin{enumerate}[nosep]
\item The gauge group at unification is Pati-Salam, not $\SU(5)$ or $\SO(10)$.
\item All three PS couplings unify: $g_4 = g_L = g_R$.
\item $\sin^2\theta_W = 3/8$ at the PS scale, running to $\approx 0.231$ at $M_Z$.
\item The Higgs is a $(\mathbf{1},\mathbf{2},\mathbf{2})$ bidoublet (Two Higgs Doublet Model).
\item Tree-level flat Higgs potential (Gauge-Higgs Unification).
\item Quartic coupling $\lambda(M_{\mathrm{PS}}) = g^2/4 \approx 0.11$.
\item Colored scalar triplets $(3,1) + (\bar{3},1)$ at $M_{\mathrm{PS}}$.
\item Proton stability: Pati-Salam preserves baryon number modulo~3, forbidding the dangerous dimension-6 proton decay operators present in $\SU(5)$~\cite{PatiSalam1974}.
\end{enumerate}

\subsection{Relation to prior work}

The signature dependence of the DeWitt supermetric on the background metric was established in generality by Schmidt~\cite{Schmidt1989}, whose formula $S = \Theta(\alpha) + s(n-s)$ (where $s$ is the number of negative eigenvalues of~$g$) yields $(6,4)$ for the Lorentzian case.
Schmidt's interest was in the variational structure of general relativity; the gauge-theoretic consequences were not explored.

Weinstein's Geometric Unity programme~\cite{Weinstein2013,Weinstein2021} independently considers the fourteen-dimensional metric bundle and identifies the Pati-Salam group via $\Spin(6,4)$; see~\cite{NguyenPolya2021} for a technical assessment.
Key differences from the present work:

\begin{enumerate}[nosep]
\item \emph{Metric:} Weinstein uses a ``chimeric'' metric; we use the standard DeWitt metric, uniquely determined by diffeomorphism invariance and ultralocality.
\item \emph{Signature:} Weinstein selects the fibre signature by an anthropic argument; we derive $(6,4)$ from an explicit computation of the DeWitt metric on a Lorentzian background.
\item \emph{Dynamics:} Weinstein introduces a novel ``shiab'' operator; we use the standard Gauss-Codazzi-Ricci decomposition of submanifold geometry (Paper~III).
\item \emph{Scope:} The present series provides explicit fermion representations (Paper~II), sign-correct gauge dynamics (Paper~III), anomaly cancellation (Paper~IV), and an analysis of the three-generation problem (Paper~V)---results not present in the Geometric Unity framework.
\end{enumerate}

The Percacci-Nesti GraviGUT programme~\cite{PercacciNesti2010} pursues a related strategy using the frame bundle rather than the metric bundle, embedding gravity and gauge fields in $\SO(3,11)$; see~\cite{KrasnovPercacci2018} for a comprehensive review of gravity-unification approaches.
Both programmes use $\GL^+(4,\R)/\SO$ as the internal space; they differ in which geometric structure (connection vs.\ metric) carries the gauge content.

\subsection{Open problems}

Several important questions remain:

\begin{enumerate}
\item \textbf{Three generations.}
The Clifford algebra $\Cl_6(\C) \cong M_8(\C)$ naturally provides one generation of fermions as $\C^8 = \mathbf{3}_{-1/3} \oplus \bar{\mathbf{3}}_{+1/3} \oplus \mathbf{1}_{-1} \oplus \mathbf{1}_{+1}$ under $\SU(3) \times \U(1)$.
Paper~V~\cite{Paper5} proposes that the three generations arise from three inequivalent complex structures on the positive-norm sector $\R^6$, and shows that anomaly constraints require $N_G \equiv 0 \pmod{3}$, which together with precision $Z$-width data forces $N_G = 3$.

\item \textbf{Hierarchy problem.}
The radiative Higgs mass is of order $(g^2/4\pi)\,M_{\mathrm{PS}} \gg 125\,\text{GeV}$.
The framework provides no natural mechanism for the required hierarchy.

\item \textbf{Fermion masses.}
The Dirac operator on~$Y$ restricted to the section has not been computed; this is needed to derive Yukawa couplings and the CKM/PMNS matrices.

\item \textbf{Quantum consistency.}
Paper~IV~\cite{Paper4} verifies that all gauge and gravitational anomalies cancel for Pati-Salam with three generations, and that the Witten $\SU(2)$ anomaly is absent.
Unitarity and UV behaviour of the fourteen-dimensional theory remain open; the submanifold approach may have better UV properties than a standard Kaluza-Klein reduction, since the fibre is non-compact and does not contribute a tower of massive modes in the same way.
Paper~VI~\cite{Paper6} shows that the Born rule, interference, and uncertainty relations emerge from the $(6,4)$ DeWitt signature when observation is modelled by a finite Markov blanket.

\item \textbf{Cosmological constant.}
The conformal dilution of the fibre curvature gives $\Lambda_{\mathrm{eff}} \sim |R_{\mathrm{fibre}}|/L_H^2$, numerically within an order of magnitude of observation, but the mechanism remains speculative.

\item \textbf{Localisation.}
The coefficient $c$ relating $G_{14}$ to $G_4$ in~\eqref{eq:action} depends on the fibre geometry.
A soldering mechanism (identifying the gauge coupling with the fibre sectional curvature) gives $\alpha \approx 0.030$ versus the observed $\alpha_{\mathrm{PS}} \approx 0.021$---a factor of $1.4$.
\end{enumerate}

\subsection{Outlook}

The central claim of this paper is modest: the Lorentzian DeWitt metric on the space of metrics over four-dimensional spacetime has signature $(6,4)$, and this fact, combined with the submanifold Gauss equation, yields the Pati-Salam gauge group with correct coupling relations.
Whether the full metric bundle framework can address the open problems listed above---and thereby constitute a viable theory of unification---remains to be seen.
The computation presented here is a necessary first step: establishing that the geometry of $\Met(X^4)$ naturally encodes the correct gauge structure of the Standard Model.


% ============================================================
% ACKNOWLEDGEMENTS
% ============================================================

\section*{Acknowledgements}

The author thanks Claude (Anthropic) for assistance with verification of the numerical computations.

% ============================================================
% REFERENCES
% ============================================================

\bibliographystyle{unsrt}
\bibliography{references}

% ============================================================
% APPENDIX
% ============================================================

\appendix

\section{Explicit matrix computation of the \texorpdfstring{$(6,4)$}{(6,4)} signature}
\label{app:matrix}

We provide the explicit $10 \times 10$ DeWitt metric matrix for verification.
With the basis ordered as $(e_{(00)}, e_{(01)}, e_{(02)}, e_{(03)}, e_{(11)}, e_{(12)}, e_{(13)}, e_{(22)}, e_{(23)}, e_{(33)})$ and background $g = \diag(-1,1,1,1)$:

The diagonal elements:
\begin{align}
\mathcal{G}_{(00)(00)} &= (g^{00})^2 - \tfrac{1}{2}(g^{00})^2 = \tfrac{1}{2}\,,\\
\mathcal{G}_{(11)(11)} &= (g^{11})^2 - \tfrac{1}{2}(g^{11})^2 = \tfrac{1}{2}\,,\\
\mathcal{G}_{(0a)(0a)} &= 2\, g^{00}g^{aa} = -2 \quad (a = 1,2,3)\,,\\
\mathcal{G}_{(ab)(ab)} &= 2\, g^{aa}g^{bb} = 2 \quad (a < b; \; a,b \in \{1,2,3\})\,.
\end{align}
The significant off-diagonal elements are:
\begin{equation}
\mathcal{G}_{(00)(aa)} = \tfrac{1}{2}\, g^{00}g^{aa} = -\tfrac{1}{2}\,, \qquad
\mathcal{G}_{(aa)(bb)} = -\tfrac{1}{2}\, g^{aa}g^{bb} = -\tfrac{1}{2} \quad (a \ne b)\,.
\end{equation}

The eigenvalues of this matrix are
\[
\{-2, -2, -2, -1, +1, +1, +1, +2, +2, +2\}\,,
\]
which after rescaling gives signature~$(6,4)$.
Independent numerical verification using two different codebases (with both $g = \diag(-1,1,1,1)$ and $g = \diag(1,1,1,-1)$) confirms this result.

The Python code accompanying this paper provides a complete, self-contained verification.

\end{document}
